% Exemplo de frame de comparação com narrativa — copiar a estrutura, não o conteúdo.
% Princípio: cada frame de módulo migrado conta uma HISTÓRIA de 3 batidas —
% (1) o que medimos, (2) o que isso significou, (3) o que decidimos e por quê —
% nunca só uma tabela solta sem a decisão que ela gerou.

\begin{frame}{RAG híbrido: Python puro vs. LangChain}
  \begin{numerobloco}{O que medimos (golden set, 49 perguntas de rota RAG)}
    \begin{center}
    \begin{tabular}{lccc}
      \toprule
      \textbf{Métrica} & \textbf{Produção (TF-IDF)} & \textbf{LangChain fiel} & \textbf{LangChain BM25+RRF} \\
      \midrule
      Recall@5    & 98,0\%  & 98,0\%  & 98,0\%  \\
      Precision@5 & 38,4\%  & 38,4\%  & \alert{49,8\%} \\
      MRR         & 0,898   & 0,898   & 0,897   \\
      \bottomrule
    \end{tabular}
    \end{center}
  \end{numerobloco}

  \pause
  \textbf{O que isso significou:} o ganho de qualidade (+11,4pp de Precision@5) vem do
  \alert{algoritmo} BM25+RRF, não do framework LangChain em si — a versão "fiel" (mesmo
  algoritmo, framework trocado) empata exatamente com a produção.

  \pause
  \begin{decisaobloco}{O que decidimos}
    Migramos mesmo assim: o roteador e o self-repair vão para LangGraph em seguida, e ter o
    RAG já em LangChain evita uma segunda migração. Usamos a versão BM25+RRF (framework +
    ganho real juntos), não a versão "fiel".
  \end{decisaobloco}
\end{frame}
